\documentclass[12pt,a4paper]{article}
\usepackage[utf8]{inputenc}
\usepackage{amsmath, amssymb, amsfonts}
\usepackage{geometry}
\geometry{margin=2.5cm}
\usepackage{hyperref}
\usepackage{pgfplots}
\usepackage{fancyhdr}
\usepackage{listing}
\usepackage{listings}
\usepackage{titling}
\usepackage{float}
\usepackage{titlesec}

\titleformat{\section}
  {\normalfont\Large\bfseries\centering}  % Format: font size, bold, centered
  {\thesection}{1em}{}
\pagestyle{fancy}
\pgfplotsset{compat=1.18}
\fancyhead{}
\fancyfoot{}
\title{Creare il profilo Linkedin}
\pretitle{\begin{center}\LARGE}
\posttitle{\par\large Career Labs\end{center}}
\author{}
\date{}
\fancyhead[L]{Gabriele Perna}          % Left side
\fancyhead[C]{Creare il profilo Linkedin}     % Center
\fancyhead[R]{2026/2027}             % Right side
\fancyfoot[C]{\thepage}           % Page number in the center of footer
\let\oldsection\section
\renewcommand{\section}{\clearpage\oldsection}

\begin{document}

\maketitle
\tableofcontents

\section{Informazioni e risorse}
\textbf{Nominativo consulente:} Giulia Verdigi.\\
\textbf{Email consulente:} giulia.verdigi@ospiti.unipi.it\\
\textbf{Instagram consulente:} gverdigi\_psicologa\\\\
\textbf{Consulenza} \textbf{online} \textbf{(gratuita) sui colloqui di lavoro:}
\href{https://alumni.unipi.it/mentoring}{https://alumni.unipi.it/mentoring}

\section{Come mai Linkedin?}
Linkedin è la piattaforma principale del mondo di lavoro e la comunità è enorme, ci sono:
\begin{itemize}
    \item 1 miliardo di membri nel mondo
    \item 67 milioni di aziende
    \item 136 mila scuole / università
\end{itemize}
Questi dati sono relativi al 2026.\\
Linkedin permette di fare \textbf{networking} e mostrare la nostra carriera.\\
In realtà si può parlare di tutto, ma sarebbe appropriato riuscire a contestualizzare i contenuti per trarne un vantaggio in ambito lavorativo.

\section{Elementi del profilo}
\subsection{Primo impatto}
Ecco gli elementi più importanti del profilo:
\begin{itemize}
    \item Foto profilo
    \item Banner (foto di sfondo)
    \item Nominativo
    \item Lingua del profilo
    \item Titolo (lavoro o specializzazione)
\end{itemize}
Questi sono gli elementi che modellano il primo impatto del profilo, se non li settiamo nel modo giusto, potremmo perdere visite dettagliate al profilo.\\
La foto profilo ideale è a mezzo busto, in quanto non si può zoommare su Linkedin.\\
La foto di sfondo deve rappresentare quello che vogliamo andare a fare nella vita, deve essere contestualizzata.\\
Il titolo dovrebbe essere specifico, ai recruiter piace chi è deciso in quello che vuole fare.
\subsection{Esperienza}
Esiste una sezione dove possiamo elencare le nostre esperienze.\\
Concettualmente, in tale sezione si includono esperienze lavorative, ma in realtà possiamo anche includere esperienze non lavorative o non retribuite che però noi stessi consideriamo importanti o simili ad un vero lavoro. Alcune esperienze non lavorative possono comunque contribuire allo sviluppo delle competenze.\\
Infatti, ogni esperienza è associata a delle competenze acquisite.
\subsection{Istruzione}
Su Linkedin, il campo "istruzione" dovrebbe includere solo i titoli conseguiti, non corsi o certificazioni, in quanto essi devono essere presentati nella loro apposita sezione.

\section{Networking}
Lo scopo di Linkedin è farsi vedere, creare connessioni utili e rimanere presente nel mondo del lavoro.\\
Un'ottima mossa strategica è seguire le aziende di nostro interesse per restare aggiornati sull'apertura di nuove posizioni, progetti e nuove strategie.\\
Inoltre, è consigliato rimanere attivi nel campo di nostro interesse per aumentare la visibilità e raffinare l'algoritmo.\\
Più connessioni abbiamo, più è alta la nostra visibilità, in quanto Linkedin mostra spesso le connessioni comuni anche non dirette, per esempio anche fino a 3 livelli di conoscenza.

\section{Cosa ho potuto migliorare nel mio profilo}
Al momento della frequentazione del corso, il sottoscritto Gabriele Perna aveva già creato un profilo Linkedin completo, per questo è stato possibile svolgere un analisi su tale profilo per individuare elementi da migliorare.\\
Ecco quali sono state le modifiche effettuate:
\begin{itemize}
    \item \textbf{Modifica dell'immagine di banner:} l'immagine precedente non era idonea in quanto rappresentava un oceano.
    \item \textbf{Zoom dell'immagine profilo o sostituzione:} aumentato lo zoom dell'immagine profilo.
    \item \textbf{Aggiunta dell'audio di pronuncia del nome}: utile per gli stranieri.
    \item \textbf{Aggiunta dei progetti nella loro sezione}: i progetti erano presentati altrove ma non inclusi nella loro sezione.
    \item \textbf{Aggiunta esperienza lavorativa come tutor di informatica}: esperienza non considerata come tale.
    \item \textbf{Rendere "open to work" visibile a tutti}: consigliato per profili junior.
    \item \textbf{Aggiunta esperienza lavorativa come modder di 3DS / PSP}: esperienza non considerata come tale.
\end{itemize}

\end{document}